% ============================================================================
%  METHOD -- drop-in section for the ASIA CCS / ACM CCS paper.
%  Describes \ourmethod{} only; the baseline collapse analysis that
%  motivates it lives in Section~\ref{sec:collapse}
%  (VLI_collapse_mechanism_en.tex) and is not repeated here.
%
%  Ported and trimmed from VLI_report.tex ("VLI and VLI*"), which is
%  correct for this content (unlike its baseline-collapse section).
%
%  Proof-status convention carried over from the Vietnamese reports:
%  [proved] = short algebraic proof, in Appendix~\ref{app:proofs}. This
%  section has no [argued] or [open] claims of its own; those live in
%  Section~\ref{sec:collapse}.
%
%  Dependencies: same theorem/macro environments as
%  VLI_preliminaries_en.tex and VLI_collapse_mechanism_en.tex.
%  \ref{sec:prelim}, \ref{sec:prelim-notation}, \ref{sec:collapse},
%  \ref{sec:eval} point at sections defined elsewhere; \label{app:proofs}
%  is bound in the appendix. \ourmethod{} is the same macro as in the
%  other files. Cites \ref{lem:tildenorm} and \ref{rem:dilemma} from
%  Section~\ref{sec:collapse} once, in Corollary~\ref{cor:no-branch}, to
%  state what \ourmethod{} avoids without re-deriving it.
%
%  NOT ported here: numerical verification tables (Section numerics of
%  VLI_report.tex) are evaluation, not method; they belong in the paper's
%  Evaluation section.
% ============================================================================

\section{Method}
\label{sec:method}

Section~\ref{sec:collapse} showed that every auxiliary-calibrated
baseline we compare against borrows, from outside the attacked window, a
quantity of \emph{absolute magnitude} that does not shrink with the
window's own training error (Lemma~\ref{lem:norm}, Proposition~\ref{prop:gap}).
This section presents \ourmethod{}, built around the opposite design
choice: a quantity of absolute magnitude never enters it at all.

\subsection{The design principle}

By Lemma~\ref{lem:identities}, $a_j = N_j r_j\bar\epsilon$: the label
information lies in the \emph{relative} structure of $G$ across classes,
while $\bar\epsilon$ -- the quantity that vanishes under training --
enters only as a common factor. An estimator that reads $G$ only through
ratios of its own coordinates therefore never needs to compare it against
anything of fixed external scale.

\begin{assumption}[Symmetric error geometry]
\label{ass:sym}
All present classes have equal mean error ($r_j=1$), every class
distributes its error mass uniformly over the others
($Q_{jk}=1/(C-1)$, $k\ne j$), at least one class is absent, and
$\bar\epsilon>0$.
\end{assumption}

Assumption~\ref{ass:sym} is false for any real model; the next result is
why it is nonetheless the right default rather than a simplification.

\begin{proposition}[Uniqueness under label exchangeability]
\label{prop:exchange}
The uniform matrix $Q_{jk}=1/(C-1)$ ($k\ne j$), $Q_{jj}=0$, is the unique
row-stochastic, zero-diagonal matrix invariant under every permutation of
class labels.
\end{proposition}

\begin{proof}[\textnormal{[proved]}]
See Appendix~\ref{app:proofs}.
\end{proof}

An attacker with no class-specific side information has no basis for
breaking symmetry among labels; Proposition~\ref{prop:exchange} says the
uniform choice is the only admissible one on these grounds, not a
convenience. Proposition~\ref{prop:weak} (below) shows this is more
robust than it looks.

\subsection{The closed form}

\begin{theorem}[Closed form]
\label{thm:closed}
Under Assumption~\ref{ass:sym},
\[
 G_j = \frac{n\bar\epsilon - Ca_j}{C-1}, \qquad
 G_{\max} = \frac{n\bar\epsilon}{C-1} \ \ (\text{at the absent class}),
 \qquad
 \boxed{\,N_j = \frac{n}{C}\Bigl(1-\frac{G_j}{G_{\max}}\Bigr)\,} ,
\]
invariant under $G \mapsto cG$ for any $c>0$.
\end{theorem}

\begin{proof}[\textnormal{[proved]}]
See Appendix~\ref{app:proofs}.
\end{proof}

$\bar\epsilon$ cancels in the last step: $G_{\max}$ is an
\emph{in-window thermometer} for $\bar\epsilon$, read from the attacked
update itself (an absent class has $a_j=0$, so its coordinate reports the
error scale directly) rather than borrowed from outside it. Where RLU,
ZLG$^+$, and LLG$^+$ borrow that scale from auxiliary data and are
punished by Proposition~\ref{prop:gap} (Section~\ref{sec:collapse}),
\ourmethod{} does not need to.

\subsection{Choosing the denominator $G^*$ in practice}
\label{sec:method-gstar}

Theorem~\ref{thm:closed} defines $G_{\max}$ as the coordinate of $G$ at
\emph{the} absent class, which Assumption~\ref{ass:sym} guarantees is
unique. Real windows can have several absent or near-absent classes, and
measurement noise can make any single coordinate of $G$ an unreliable
denominator on its own: if the coordinate used as $G^*$ is inflated by a
factor $(1+b)$ over its ideal value, every recovered count shrinks toward
the uniform prior by the same factor,
$\widehat N_j = \tfrac nC + \lambda\bigl(N_j - \tfrac nC\bigr)$ with
$\lambda = 1/(1+b)$ -- so which coordinate, or combination of
coordinates, is used as $G^*$ is a practical choice with a direct,
measurable effect on recovery, not a detail. We implement three:

\begin{table}[t]
\centering
\caption{Choices of $G^*$ in $N_j = \tfrac nC(1-G_j/G^*)$. $S^+ = \{j :
G_j>0\}$ denotes the positive coordinates of $G$.}
\label{tab:gstar}
\small
\begin{tabular}{@{}llp{3.1cm}@{}}
\toprule
Variant & $G^*$ & Exact under Assumption~\ref{ass:sym} \emph{iff} \\
\midrule
\texttt{max}    & $\max_j G_j$                 & always (by construction) \\
\texttt{mean+}  & $\mathrm{mean}_{j\in S^+} G_j$ & $|S^+|=1$ \\
\texttt{med+}   & $\mathrm{median}_{j\in S^+} G_j$ & $|S^+|=1$ \\
\bottomrule
\end{tabular}
\end{table}

\texttt{max} is the literal reading of Theorem~\ref{thm:closed} and is
exact whenever the theorem's hypotheses hold, but a single coordinate is
also the most exposed to noise inflating $b$. \texttt{mean+} and
\texttt{med+} average over every positive coordinate instead of trusting
one, which costs exactness when more than one class is genuinely absent
($|S^+|>1$, where the theorem no longer pins all of them to the same
value) but empirically reduces $b$ when the nominal maximum is itself an
outlier -- \texttt{med+} in particular discards exactly the inflation
that would otherwise dominate \texttt{mean+}'s average. \texttt{med+} is
the $q=0.5$ case of the one-parameter family
$G^*_q = \mathrm{quantile}_q(\{G_j\}_{j\in S^+})$ that the implementation
also supports; we report \texttt{max}, \texttt{mean+}, and \texttt{med+}
as the three fixed points of this family worth naming, and leave sweeping
$q$ to Section~\ref{sec:eval}.

\begin{corollary}[\ourmethod{} does not inherit the weight-branch dilemma]
\label{cor:no-branch}
\ourmethod{} reads $G$ directly (Table~\ref{tab:baselines}), never
$\widetilde G$. Consequently Lemma~\ref{lem:tildenorm}'s residual term,
and the dilemma of Remark~\ref{rem:dilemma} (Section~\ref{sec:collapse}),
do not arise for \ourmethod{} regardless of how $w_i$ behaves.
\end{corollary}

\begin{proposition}[Scale invariance and structural validity]
\label{prop:invariance}
The map $\Psi: G \mapsto N$ of Theorem~\ref{thm:closed} satisfies
$\Psi(cG)=\Psi(G)$ for every $c>0$; its output is nonnegative and sums to
$n$ by construction (no clipping or renormalization step is required),
and under Assumption~\ref{ass:sym} it reproduces $N_j - N_k$ exactly in
its coordinate differences.
\end{proposition}

\begin{proof}[\textnormal{[proved]}]
See Appendix~\ref{app:proofs}.
\end{proof}

No external quantity of absolute magnitude enters $\Psi$ at all, so
Theorem~\ref{thm:collapse} (Section~\ref{sec:collapse}) cannot describe
it: Proposition~\ref{prop:invariance}'s first clause is the reason, not
merely a consequence. Needing no clipping or renormalization step also
means \ourmethod{} never exercises Proposition~\ref{prop:lip}'s
non-Lipschitz map; its repair, when one is needed, is the non-expansive
projection $\Pi_\simplex$ instead.

\begin{proposition}[A weaker sufficient condition]
\label{prop:weak}
Theorem~\ref{thm:closed} holds exactly whenever $r_j=1$ on present
classes, at least one class is absent, and the incoming error mass
satisfies $(Q^Ta)_j = (\sum_k a_k - a_j)/(C-1)$ for every $j$ -- a
constraint on the \emph{aggregate} inflow to each class, strictly weaker
than requiring $Q$ itself to be uniform.
\end{proposition}

\begin{proof}[\textnormal{[proved]}]
See Appendix~\ref{app:proofs}.
\end{proof}

\subsection{Oracle instruments for the ablation}
\label{sec:method-oracle}

Assumption~\ref{ass:sym} does not hold for any real model. Therefore, the
useful question is how much it costs. The assumption makes two
idealizations. First, all present classes have the same difficulty, so
\(r_j \equiv 1\). Second, confusion is uniform, so
\(Q_{jk} = 1/(C-1)\). ZLI also has to choose a denominator \(G^*\), as
described in Section~\ref{sec:method-gstar}. This gives three separate
sources of error.

Each source can be replaced by its true value, computed from a labelled
window. An adversary in Section~\ref{sec:threat} cannot do this. We use it
only as an experimental instrument. It lets us attribute the remaining
error of ZLI to a specific cause.

The next theorem gives the ceiling. If both \(r\) and \(Q\) are known,
recovery is exact.

\begin{theorem}[Recovery with known difficulty and confusion geometry]
\label{thm:general}
If the per-class difficulty \(r\) and the confusion matrix \(Q\)
(Section~\ref{sec:prelim-notation}) were known exactly, the unique
nonnegative solution of \(G=-(I-Q^{\top})a\) with a zero coordinate at
the known absent class recovers the true counts exactly.
\end{theorem}

\begin{proof}[\textnormal{[proved]}]
See Appendix~\ref{app:proofs}.
\end{proof}

Corollary~\ref{cor:kappa} gives a prediction for one arm. It bounds the
loss when \(Q\) is known but \(r\) is wrongly assumed to be uniform. We
can compare the measured result of that arm against this bound.

\begin{corollary}[Cost of assuming \(r\equiv1\) under known \(Q\)]
\label{cor:kappa}
If \(Q\) is known and exact but \(r\) is wrongly taken to be uniform, the
instance-recovery rate is at least \(1/\kappa\), where
\(\kappa=r_{\max}/r_{\min}\) over present classes.
\end{corollary}

\begin{proof}[\textnormal{[proved]}]
See Appendix~\ref{app:proofs}.
\end{proof}

Table~\ref{tab:oracle} lists the arms.

\begin{table}[t]
\centering
\caption{Oracle arms of the ablation (Section~\ref{sec:eval}). Each arm
replaces one idealization of ZLI with its true value from the labelled
window.}
\label{tab:oracle}
\footnotesize
\setlength{\tabcolsep}{3pt}
\begin{tabular}{@{}p{1.3cm} p{1.9cm} p{3.1cm}@{}}
\toprule
Arm & Replaces & Expected effect \\
\midrule
oracle-\(Q\) & true \(Q\); \(r\equiv1\) kept & irec \(\ge 1/\kappa\) (Cor.~\ref{cor:kappa}) \\
oracle-\(r\) & true \(r\); uniform \(Q\) kept & not bounded; measured \\
oracle-\(G^*\) & true mean over absent classes & isolates the denominator \\
\midrule
oracle-all & true \(r\) and \(Q\) & exact (Thm.~\ref{thm:general}) \\
\bottomrule
\end{tabular}
\end{table}

The arms answer two different questions. The \(r\) and \(Q\) arms measure
the cost of the symmetry assumption. The \(G^*\) arm measures the cost of
estimating the denominator from the signal. The arms interact, so they are
not additive. We therefore report each arm separately, and use
\textsc{oracle-all} as the final ceiling. The choice between
\texttt{max} and \texttt{mean+} affects only the third arm. Thus,
\textsc{oracle-}\(G^*\) is the ceiling for that choice.
